\chapter[Discusión y conclusiones]{Discusión y conclusiones\label{cap:discusion}}

\section{Discusión}

\label{sec:discussion}

El desarrollo presentado en esta tesis debe interpretarse como la integración de
un flujo observacional completo, y no únicamente como la suma de modelos de
aprendizaje profundo evaluados de forma aislada. A lo largo del trabajo se
articulan decisiones instrumentales, criterios de construcción de datos,
modelos de visión por computadora, herramientas predictivas y mecanismos de
despliegue local y web, siguiendo la arquitectura metodológica formulada en el
Capítulo~\ref{cap:consideraciones}, implementada en el Capítulo~\ref{cap:software}
y validada en el Capítulo~\ref{cap:resultados}. Esta perspectiva resulta importante
porque el problema de análisis automático de manchas solares no se limita a
reconocer patrones en una imagen: exige capturar observaciones de manera
reproducible, asociarlas con metadatos físicos, preservar trazabilidad, ejecutar
inferencia en condiciones de campo y ofrecer resultados interpretables para el
usuario.

Desde el punto de vista instrumental, la selección de una cadena óptica basada en
observación de luz blanca, cámara monocromática y atenuación mediante cuña de
Herschel responde a un compromiso entre seguridad, resolución, portabilidad y
viabilidad operativa. La ZWO ASI174MM Mini ofrece un sensor de área suficiente,
píxeles relativamente grandes, obturador global y compatibilidad práctica con
entornos Linux, rasgos adecuados para una estación de observación local. La
Raspberry Pi 5, por su parte, permite concentrar adquisición, preprocesamiento,
conectividad y servicios locales en una unidad compacta. Esta combinación,
discutida en el diseño instrumental y validada físicamente en el
Capítulo~\ref{cap:resultados}, no pretende competir con observatorios solares
profesionales, sino demostrar que un observatorio académico o local puede
incorporar una capa de asistencia computacional moderna sin requerir
infraestructura de gran escala \citep{zwoASI174Manual2025,luntWhiteLightWedges2026,raspberryPi5Docs2026}.

La discusión de hardware también revela una restricción central del sistema: la
calidad de la inferencia depende de la estabilidad de la adquisición. Condiciones
como el \textit{seeing}, la nubosidad, el enfoque, la exposición y la ganancia de
la cámara condicionan la nitidez de umbra y penumbra, y por tanto afectan tanto
la detección de grupos como la clasificación morfológica. Por ello, el
preprocesamiento local, el control de parámetros de captura y la verificación de
estado mediante panel local no son accesorios secundarios, sino componentes
necesarios para sostener la calidad de los datos que alimentan los modelos. Esta
lectura es coherente con las limitaciones de validación señaladas en el
Capítulo~\ref{cap:resultados}, donde se reconoce que las campañas observacionales
están condicionadas por clima, ventana horaria solar y estabilidad instrumental.

En relación con los datos de entrenamiento, una de las decisiones metodológicas
más relevantes fue utilizar imágenes SDO/HMI acopladas con reportes NOAA SRS para
construir conjuntos etiquetados, como se detalla en la
Sección~\ref{sec:data-software}. Esta elección se justifica por la disponibilidad
temporal, la regularidad instrumental y la existencia de metadatos taxonómicos
asociados \citep{scherrer2012hmi,noaaSRS}. Sin embargo, también introduce
una tensión metodológica: el sistema
final está orientado a imágenes de luz blanca obtenidas localmente, mientras que
el entrenamiento se apoya en una modalidad espacial fotosférica distinta. La
similitud morfológica entre HMI y luz blanca permite sostener esta aproximación
como punto de partida razonable, pero no elimina la necesidad de validar la
transferencia entre instrumentos. En consecuencia, los resultados obtenidos deben
leerse como una validación funcional y metodológica del flujo, no como una
caracterización definitiva del rendimiento sobre todos los escenarios de
observación terrestre.

La separación entre agrupamiento y clasificación constituye otro aspecto central
de la propuesta. Esta estrategia en cascada fue formalizada en la
Sección~\ref{sec:arquitectura-dl-cascada} e implementada en la
Sección~\ref{sec:training}. El detector YOLO26n opera sobre imágenes de disco
completo y tiene como objetivo localizar grupos solares, mientras que el
clasificador basado en ConvNeXtV2 Atto trabaja sobre recortes normalizados y
estima las componentes Zurich--McIntosh \citep{ultralyticsYOLO26,woo2023convnextv2}.
La variante Atto fue seleccionada para preservar la integridad de la red durante
la transformación desde PyTorch hacia ONNX y para sostener un despliegue ligero.
Sobre este backbone se inyectan etapas complementarias basadas en MLP que
implementan el aprendizaje jerárquico entre las salidas $Z$, $P$ y $C$.
Esta división reduce la complejidad de cada tarea y se alinea
con la lógica observacional: primero se identifica la unidad física relevante,
el grupo solar, y luego se describe su morfología. Los resultados de detección,
con precisión del orden del 83\%, indican que la primera etapa alcanza un punto
de operación suficiente para alimentar la cadena posterior. No obstante, los
errores de localización siguen siendo críticos, ya que un recorte incompleto,
descentrado o contaminado puede degradar la clasificación aunque el clasificador
sea correcto en condiciones ideales.

La clasificación jerárquica Zurich--McIntosh muestra un comportamiento
competitivo frente a referencias recientes. Según la comparación de la
Tabla~\ref{tab:comparativa-baselines}, el F1 promedio de 0.66 se aproxima a los
mejores valores considerados, aun sin recurrir a compactación artificial de
clases ni a ensambles de modelos
\citep{palladino2022sunspot,sayez2023sunscc,deng2025mcintosh}. Esta
comparación es relevante porque la
propuesta conserva la estructura taxonómica completa y mantiene salidas
separadas para $Z$, $P$ y $C$. Desde el punto de vista científico, esta decisión
favorece la interpretabilidad: el sistema no solo entrega una etiqueta final,
sino que permite observar cómo se distribuye la probabilidad entre alternativas
compatibles. El árbol bayesiano y el top 12 de recomendaciones, presentados en
las Figuras~\ref{fig:resultados-arbol-top12-zpc} y~\ref{fig:resultados-top12-zpc},
implementan la normalización bayesiana descrita en la
Sección~\ref{sec:classification}; por ello, la salida puede funcionar como una
herramienta de asistencia, ofreciendo una clase principal y alternativas
plausibles cuando existen configuraciones cercanas.

El desbalance de clases sigue siendo una de las principales limitaciones. La
taxonomía Zurich--McIntosh presenta una distribución de cola larga: pocas
combinaciones concentran gran parte de los ejemplos y muchas clases aparecen de
forma escasa. Esta lectura se apoya en la distribución de DB2 y en las matrices
de confusión por componente, sin requerir un análisis cualitativo separado de
casos particulares. La estructura observada afecta especialmente a
configuraciones raras o morfológicamente ambiguas. Además, las fronteras entre
ciertas clases no siempre son nítidas incluso para observadores humanos, lo que
introduce ruido de etiqueta difícil de eliminar por completo. En este contexto,
la estrategia jerárquica y la máscara de combinaciones válidas ayudan a
restringir el espacio de salida, pero no sustituyen la necesidad de ampliar y
depurar los datos.

Los modelos predictivos incorporados cumplen una función complementaria. El
modelo PVL con filtro de Kalman ofrece una trayectoria escalar suavizada del
número de manchas solares, mientras que el modelo DMD introduce una visión
espacio-temporal de la migración latitudinal, de acuerdo con la formalización de
la Sección~\ref{sec:forecast}. Los errores MAPE reportados en la
Tabla~\ref{tab:ssn-metrics-comparison} indican que el enfoque PVLK resulta más
preciso en términos escalares, pero las métricas de la Tabla~\ref{tab:sporer-performance}
muestran que el modelo DMD aporta información física adicional al representar la
distribución latitudinal y su desplazamiento hacia el ecuador
\citep{kalman1960newapproach,podladchikova2012kalman,schmid2010dmd,kutz2016dmd}.
Esta diferencia ilustra un aspecto importante
del sistema: no todos los módulos deben optimizar la misma métrica. En algunos
casos, una representación menos precisa en un índice global puede ser más rica
para interpretación científica si conserva estructura espacial o temporal.

La plataforma web y el despliegue ONNX completan la dimensión operativa del
trabajo, desarrollada en la Sección~\ref{sec:web}. Al convertir los modelos a
ONNX y ejecutarlos mediante mecanismos compatibles con navegador o dispositivos
locales, la tesis reduce la dependencia de un servidor central de inferencia. El
peso conjunto de los modelos ONNX, cercano a 24~MB, resulta especialmente
conveniente para descarga y actualización en entornos de conectividad limitada.
Esto mejora portabilidad, disminuye la latencia y permite que parte del
procesamiento ocurra cerca del lugar de captura.
La integración con base de datos y almacenamiento en la nube añade trazabilidad:
las imágenes, metadatos, recortes e inferencias pueden consultarse posteriormente,
compararse y auditarse. Esta trazabilidad es especialmente importante en un
contexto científico, donde cada resultado debe poder asociarse con su fecha,
origen instrumental y etapa de procesamiento.

La validación integral confirma que el prototipo no quedó limitado a una prueba
de laboratorio. Las evidencias de integración física, operación local, acceso web,
procesamiento ONNX y visualización de resultados, reunidas en el
Capítulo~\ref{cap:resultados}, muestran que la arquitectura puede ejecutarse como
sistema funcional. Sin embargo, la validación de campo aún
es acotada. Falta evaluar el rendimiento durante campañas observacionales más
largas, bajo diferentes condiciones atmosféricas, con variaciones de enfoque y
exposición, y con una muestra mayor de imágenes capturadas directamente por el
prototipo. Esta etapa será decisiva para cuantificar la robustez real del sistema
en el dominio de aplicación final.

En síntesis, la principal contribución de la tesis no radica únicamente en haber
entrenado un detector o un clasificador, sino en haber construido una arquitectura
coherente que conecta adquisición, datos, modelos, inferencia, visualización y
análisis. El sistema conserva vínculos explícitos con la interpretación solar:
localiza grupos, predice componentes de una taxonomía observacional, restringe
combinaciones inválidas, modela tendencias temporales y mantiene trazabilidad de
los productos generados. Estas características lo convierten en una base sólida
para futuras extensiones hacia observación solar asistida en entornos académicos
y observatorios locales.

\section{Conclusiones}

\label{sec:conclusion}

Esta tesis presentó un prototipo integral para la detección, clasificación y
análisis predictivo de manchas solares mediante visión computacional e
inteligencia artificial. El sistema propuesto combina una arquitectura de
adquisición local, una cadena de datos basada en imágenes SDO/HMI y registros
NOAA SRS, modelos de aprendizaje profundo para agrupamiento y clasificación, y
herramientas de análisis temporal orientadas al número de manchas solares y a la
distribución latitudinal de la actividad. La propuesta queda articulada entre el
diseño metodológico del Capítulo~\ref{cap:consideraciones}, la implementación del
Capítulo~\ref{cap:software} y la validación del Capítulo~\ref{cap:resultados}.

En relación con el diseño instrumental, se definió una configuración viable para
observación solar de luz blanca basada en cámara monocromática, cuña de Herschel,
Raspberry Pi 5, pantalla local de estado, conectividad y carcasa de protección,
coherente con los criterios de selección técnica discutidos en el
Capítulo~\ref{cap:consideraciones}.
Esta configuración permite articular captura, preprocesamiento y sincronización
de datos en un dispositivo compacto. La selección de componentes respondió a
criterios de compatibilidad, portabilidad, seguridad óptica y capacidad de
operación local.

Respecto a la construcción de datos, se estableció un flujo de acoplamiento entre
imágenes HMI y reportes NOAA SRS para generar etiquetas útiles para entrenamiento
supervisado. Este procedimiento permitió construir una base para detección de
grupos solares y otra para clasificación Zurich--McIntosh. La tesis reconoce que
esta solución responde a una restricción práctica de disponibilidad de datos:
aunque el objetivo aplicado son imágenes de luz blanca, el entrenamiento se apoyó
en HMI por su cobertura temporal, estabilidad y relación con registros
catalográficos.

En la etapa de detección, el modelo YOLO26n alcanzó un desempeño adecuado para
localizar regiones activas sobre imágenes de disco completo, según la evaluación
resumida en la Sección~\ref{sec:results}. Su función dentro
del sistema no es segmentar detalles finos, sino producir candidatos espaciales
consistentes para el clasificador. Esta decisión resultó apropiada para una
arquitectura en cascada, donde cada módulo resuelve una parte específica del
problema observacional.

En la etapa de clasificación, la arquitectura basada en ConvNeXtV2 Atto y salidas
jerárquicas para $Z$, $P$ y $C$ mostró un rendimiento competitivo frente a
trabajos recientes y mantuvo correspondencia con la formulación probabilística de
la Sección~\ref{sec:classification}. La incorporación de etapas MLP
complementarias permitió inyectar la dependencia jerárquica entre componentes sin
comprometer la exportación íntegra del modelo desde PyTorch hacia ONNX. El F1
promedio de 0.66 confirma que la propuesta puede
aproximarse al estado del arte manteniendo interpretabilidad y sin recurrir a
compactación artificial de clases. La incorporación de restricciones
probabilísticas y combinaciones taxonómicas válidas permite transformar la salida
del modelo en una recomendación más útil para revisión humana.

En cuanto al análisis predictivo, los modelos PVLK y DMD aportan perspectivas
complementarias, consistentes con la formulación de la Sección~\ref{sec:forecast}. PVLK ofrece una estimación escalar más precisa del número de
manchas solares suavizado, mientras que DMD permite reconstruir la dinámica
latitudinal asociada al diagrama de Spörer. Estos módulos no deben interpretarse
como sistemas operacionales de pronóstico espacial, sino como herramientas de
apoyo exploratorio integradas dentro de la plataforma.

La conversión de modelos a ONNX y su despliegue en una plataforma web muestran,
como se describe en la Sección~\ref{sec:web}, que el sistema puede trascender el
entorno de entrenamiento. El tamaño conjunto cercano a 24~MB hace viable la
distribución de los modelos incluso en escenarios de conectividad limitada. La
ejecución local o en cliente reduce dependencias, mejora la portabilidad y
fortalece la autonomía del prototipo. La integración con base de datos y
almacenamiento remoto permite
conservar trazabilidad completa entre imagen, metadatos, detección,
clasificación e historial de observación.

Por tanto, se concluye que la separación modular entre adquisición,
preprocesamiento, detección, clasificación, predicción y visualización constituye
una estrategia adecuada para el análisis asistido de manchas solares. El trabajo
cumple su objetivo principal al demostrar la viabilidad técnica de un sistema
integrado para observación solar local, con capacidad de procesamiento
inferencial y soporte interpretativo. Sus limitaciones principales se relacionan
con la transferencia HMI--luz blanca, el desbalance de clases, el ruido de
etiquetado, la validación de campo aún acotada por condiciones climáticas y
ventanas horarias de observación, y la necesidad de ampliar la base de datos con
observaciones propias.

\section{Trabajos futuros}

Los resultados obtenidos abren varias líneas de continuidad. La primera y más
importante es la construcción de una base de imágenes de luz blanca etiquetada a
partir de observaciones locales. Esta base debería incluir imágenes capturadas en
distintas condiciones atmosféricas, con diferentes configuraciones de exposición,
ganancia y enfoque, y con anotaciones verificadas por expertos o mediante revisión
asistida. Con ello sería posible medir de forma directa la transferencia entre el
entrenamiento HMI--SRS y el dominio instrumental final de luz blanca.

Una segunda línea consiste en mejorar la curación de etiquetas. El uso de
Roboflow para revisar cajas y corregir anotaciones automáticas puede ampliarse
mediante protocolos de doble revisión, medición de acuerdo entre observadores y
registro explícito de incertidumbre en etiquetas ambiguas. En lugar de tratar la
etiqueta como una verdad única, futuros trabajos podrían representar la
incertidumbre morfológica como una distribución de clases, especialmente en
fronteras difíciles entre configuraciones Zurich--McIntosh.

También será necesario profundizar el manejo del desbalance. Entre las opciones
posibles se encuentran muestreo estratificado, pérdidas focales, ponderación por
frecuencia efectiva, generación controlada de aumentos para clases raras y
esquemas de aprendizaje métrico. Otra alternativa es definir niveles de operación
según el uso: una clasificación completa para análisis detallado y una versión
compactada para escenarios donde se priorice robustez sobre granularidad.

En el plano del modelado visual, futuras versiones podrían comparar ConvNeXtV2
con arquitecturas más recientes, incluyendo modelos híbridos CNN--Transformer,
redes ligeras optimizadas para dispositivos de borde y estrategias de aprendizaje
auto-supervisado. También sería valioso incorporar información multibanda o
magnetográfica cuando esté disponible, ya que la clasificación morfológica en luz
blanca puede beneficiarse de señales complementarias sobre complejidad magnética.

Para el detector, una extensión natural consiste en evaluar modelos exportados a
ONNX directamente sobre Raspberry Pi 5 y en navegador bajo diferentes cargas de
trabajo. Deberían medirse latencia, uso de memoria, consumo energético,
estabilidad térmica y comportamiento durante sesiones prolongadas. Esta
caracterización permitiría transformar la validación funcional actual en una
validación cuantitativa del desempeño embebido.

En cuanto a los modelos predictivos, el bloque PVLK puede extenderse incorporando
actualización en línea con observaciones recientes, estimación explícita de
intervalos de confianza y comparación con modelos estadísticos alternativos. El
bloque DMD puede enriquecerse con variantes multirresolución, DMD con control,
DMD kernelizado o modelos no lineales que preserven la estructura latitudinal del
diagrama de mariposa. Una línea particularmente interesante sería combinar la
salida escalar del SSN con la dinámica latitudinal para generar indicadores
conjuntos de actividad.

Desde la perspectiva de plataforma, futuros trabajos deberían fortalecer la
interfaz de revisión humana. Esto incluye mostrar incertidumbre, alternativas top
$k$, historial de inferencias, comparación entre observaciones de distintas
fechas y herramientas para corregir manualmente resultados. Cada corrección
podría reingresar al sistema como dato validado, habilitando un ciclo de mejora
continua mediante aprendizaje activo.

Finalmente, se propone ampliar la validación institucional del prototipo mediante
campañas coordinadas con observatorios locales. Estas campañas permitirían
comparar observaciones del sistema con registros profesionales, evaluar
reproducibilidad, medir sensibilidad ante condiciones atmosféricas y consolidar
un protocolo de uso. Con estas extensiones, el prototipo podría evolucionar desde
una prueba funcional de tesis hacia una herramienta estable de apoyo para
observación solar, formación académica y generación de datos heliográficos
locales.